\chapter{Dağılım Tabloları}
\label{app:distribution-tables}
Bu ek, kitap boyunca kullanılan temel dağılım tablolarını ayrı sayfalarda ve
hiyerarşik sırada verir. Nicelik gösterimi birikimli olasılığa dayanır:
$z_q$, $t_{q,\nu}$ ve $\chi^2_{q,\nu}$ değerlerinin solunda $q$ olasılık
bulunur. Sağ kuyrukta $\alpha$ olasılık bırakmak için $q=1-\alpha$ alınır.

\begin{prismbox}[PrismLight]{Tabloyu kullanma sırası}
\begin{enumerate}
  \item Dağılımı ve serbestlik derecesini belirleyin.
  \item Sorunun tek kuyruklu, çift kuyruklu veya güven aralığı olduğunu saptayın.
  \item Sağ kuyruk olasılığını birikimli olasılığa $q=1-\alpha$ ile dönüştürün.
  \item Satır ve sütunun kesişimindeki kritik değeri okuyun.
  \item Tablodaki yuvarlama nedeniyle sınırda kalan sonuçlarda yazılımın tam
  \pval-değerini veya niceliğini kullanın.
\end{enumerate}
\end{prismbox}

Tablolarda $q$, kritik değerin solunda kalan olasılığı gösterir. Örneğin
yüzde 95 çift yönlü güven aralığında her kuyrukta $0{,}025$ kaldığından
$q=0{,}975$ sütunu kullanılır. Tek yönlü $\alpha=0{,}05$ testinde ise
$q=0{,}95$ alınır. Serbestlik derecesi tabloda yoksa en yakın küçük dereceyi
kullanmak daha koruyucu bir yaklaşık sonuç verir; yayımlanacak analizlerde
doğrudan yazılım niceliği tercih edilmelidir.

\begin{table}[H]
\centering
\caption{Sık kullanılan çift yönlü normal kritik değerler.}
\label{tab:common-z-critical}
\begin{tabular}{rrrrrr}
\toprule
Güven düzeyi & Yüzde 80 & Yüzde 90 & Yüzde 95 & Yüzde 98 & Yüzde 99 \\
\midrule
$z_{1-\alpha/2}$ & 1,282 & 1,645 & 1,960 & 2,326 & 2,576 \\
\bottomrule
\end{tabular}
\end{table}


\clearpage
\section{Standart Normal Dağılım Tablosu}
\label{sec:normal-table}
\textit{Kullanım:} Hücreler $\Phi(z)=\P(Z\le z)$ değerini verir. Negatif değerler için simetri kullanılır: $\Phi(-z)=1-\Phi(z)$.

\begin{center}
\scriptsize
\setlength{\tabcolsep}{3.2pt}
\resizebox{\textwidth}{!}{%
\begin{tabular}{c*{10}{r}}
\toprule
$z$ & 0,00 & 0,01 & 0,02 & 0,03 & 0,04 & 0,05 & 0,06 & 0,07 & 0,08 & 0,09 \\
\midrule
0,0 & 0,5000 & 0,5040 & 0,5080 & 0,5120 & 0,5160 & 0,5199 & 0,5239 & 0,5279 & 0,5319 & 0,5359 \\
0,1 & 0,5398 & 0,5438 & 0,5478 & 0,5517 & 0,5557 & 0,5596 & 0,5636 & 0,5675 & 0,5714 & 0,5753 \\
0,2 & 0,5793 & 0,5832 & 0,5871 & 0,5910 & 0,5948 & 0,5987 & 0,6026 & 0,6064 & 0,6103 & 0,6141 \\
0,3 & 0,6179 & 0,6217 & 0,6255 & 0,6293 & 0,6331 & 0,6368 & 0,6406 & 0,6443 & 0,6480 & 0,6517 \\
0,4 & 0,6554 & 0,6591 & 0,6628 & 0,6664 & 0,6700 & 0,6736 & 0,6772 & 0,6808 & 0,6844 & 0,6879 \\
0,5 & 0,6915 & 0,6950 & 0,6985 & 0,7019 & 0,7054 & 0,7088 & 0,7123 & 0,7157 & 0,7190 & 0,7224 \\
0,6 & 0,7257 & 0,7291 & 0,7324 & 0,7357 & 0,7389 & 0,7422 & 0,7454 & 0,7486 & 0,7517 & 0,7549 \\
0,7 & 0,7580 & 0,7611 & 0,7642 & 0,7673 & 0,7704 & 0,7734 & 0,7764 & 0,7794 & 0,7823 & 0,7852 \\
0,8 & 0,7881 & 0,7910 & 0,7939 & 0,7967 & 0,7995 & 0,8023 & 0,8051 & 0,8078 & 0,8106 & 0,8133 \\
0,9 & 0,8159 & 0,8186 & 0,8212 & 0,8238 & 0,8264 & 0,8289 & 0,8315 & 0,8340 & 0,8365 & 0,8389 \\
1,0 & 0,8413 & 0,8438 & 0,8461 & 0,8485 & 0,8508 & 0,8531 & 0,8554 & 0,8577 & 0,8599 & 0,8621 \\
1,1 & 0,8643 & 0,8665 & 0,8686 & 0,8708 & 0,8729 & 0,8749 & 0,8770 & 0,8790 & 0,8810 & 0,8830 \\
1,2 & 0,8849 & 0,8869 & 0,8888 & 0,8907 & 0,8925 & 0,8944 & 0,8962 & 0,8980 & 0,8997 & 0,9015 \\
1,3 & 0,9032 & 0,9049 & 0,9066 & 0,9082 & 0,9099 & 0,9115 & 0,9131 & 0,9147 & 0,9162 & 0,9177 \\
1,4 & 0,9192 & 0,9207 & 0,9222 & 0,9236 & 0,9251 & 0,9265 & 0,9279 & 0,9292 & 0,9306 & 0,9319 \\
1,5 & 0,9332 & 0,9345 & 0,9357 & 0,9370 & 0,9382 & 0,9394 & 0,9406 & 0,9418 & 0,9429 & 0,9441 \\
1,6 & 0,9452 & 0,9463 & 0,9474 & 0,9484 & 0,9495 & 0,9505 & 0,9515 & 0,9525 & 0,9535 & 0,9545 \\
1,7 & 0,9554 & 0,9564 & 0,9573 & 0,9582 & 0,9591 & 0,9599 & 0,9608 & 0,9616 & 0,9625 & 0,9633 \\
1,8 & 0,9641 & 0,9649 & 0,9656 & 0,9664 & 0,9671 & 0,9678 & 0,9686 & 0,9693 & 0,9699 & 0,9706 \\
1,9 & 0,9713 & 0,9719 & 0,9726 & 0,9732 & 0,9738 & 0,9744 & 0,9750 & 0,9756 & 0,9761 & 0,9767 \\
2,0 & 0,9772 & 0,9778 & 0,9783 & 0,9788 & 0,9793 & 0,9798 & 0,9803 & 0,9808 & 0,9812 & 0,9817 \\
2,1 & 0,9821 & 0,9826 & 0,9830 & 0,9834 & 0,9838 & 0,9842 & 0,9846 & 0,9850 & 0,9854 & 0,9857 \\
2,2 & 0,9861 & 0,9864 & 0,9868 & 0,9871 & 0,9875 & 0,9878 & 0,9881 & 0,9884 & 0,9887 & 0,9890 \\
2,3 & 0,9893 & 0,9896 & 0,9898 & 0,9901 & 0,9904 & 0,9906 & 0,9909 & 0,9911 & 0,9913 & 0,9916 \\
2,4 & 0,9918 & 0,9920 & 0,9922 & 0,9925 & 0,9927 & 0,9929 & 0,9931 & 0,9932 & 0,9934 & 0,9936 \\
2,5 & 0,9938 & 0,9940 & 0,9941 & 0,9943 & 0,9945 & 0,9946 & 0,9948 & 0,9949 & 0,9951 & 0,9952 \\
2,6 & 0,9953 & 0,9955 & 0,9956 & 0,9957 & 0,9959 & 0,9960 & 0,9961 & 0,9962 & 0,9963 & 0,9964 \\
2,7 & 0,9965 & 0,9966 & 0,9967 & 0,9968 & 0,9969 & 0,9970 & 0,9971 & 0,9972 & 0,9973 & 0,9974 \\
2,8 & 0,9974 & 0,9975 & 0,9976 & 0,9977 & 0,9977 & 0,9978 & 0,9979 & 0,9979 & 0,9980 & 0,9981 \\
2,9 & 0,9981 & 0,9982 & 0,9982 & 0,9983 & 0,9984 & 0,9984 & 0,9985 & 0,9985 & 0,9986 & 0,9986 \\
3,0 & 0,9987 & 0,9987 & 0,9987 & 0,9988 & 0,9988 & 0,9989 & 0,9989 & 0,9989 & 0,9990 & 0,9990 \\

\bottomrule
\end{tabular}
}
\end{center}

\textit{Okuma örneği:} $z=1{,}96$ için 1,9 satırı ile 0,06 sütununun
kesişiminde $\Phi(1{,}96)=0{,}9750$ bulunur. Sağ kuyruk olasılığı
$1-0{,}9750=0{,}0250$, çift yönlü olasılık ise yaklaşık $0{,}0500$'dır.


\clearpage
\section[Student t Dağılımı Kritik Değerleri]{Student $t$ Dağılımı Kritik Değerleri}
\label{sec:t-table}
\textit{Kullanım:} Hücreler $\P(T_\nu\leq t_{q,\nu})=q$ olacak birikimli
olasılık nicelikleridir. İki taraflı $100(1-\alpha)\%$ güven aralığında
$q=1-\alpha/2$ alınır.

\begin{center}
\scriptsize
\renewcommand{\arraystretch}{0.84}
\setlength{\tabcolsep}{5pt}
\resizebox{\textwidth}{!}{%
\begin{tabular}{c*{5}{r}}
\toprule
$\nu$ & $q=0{,}90$ & $q=0{,}95$ & $q=0{,}975$ & $q=0{,}99$ & $q=0{,}995$ \\
\midrule
1 & 3,078 & 6,314 & 12,706 & 31,821 & 63,657 \\
2 & 1,886 & 2,920 & 4,303 & 6,965 & 9,925 \\
3 & 1,638 & 2,353 & 3,182 & 4,541 & 5,841 \\
4 & 1,533 & 2,132 & 2,776 & 3,747 & 4,604 \\
5 & 1,476 & 2,015 & 2,571 & 3,365 & 4,032 \\
6 & 1,440 & 1,943 & 2,447 & 3,143 & 3,707 \\
7 & 1,415 & 1,895 & 2,365 & 2,998 & 3,499 \\
8 & 1,397 & 1,860 & 2,306 & 2,896 & 3,355 \\
9 & 1,383 & 1,833 & 2,262 & 2,821 & 3,250 \\
10 & 1,372 & 1,812 & 2,228 & 2,764 & 3,169 \\
11 & 1,363 & 1,796 & 2,201 & 2,718 & 3,106 \\
12 & 1,356 & 1,782 & 2,179 & 2,681 & 3,055 \\
13 & 1,350 & 1,771 & 2,160 & 2,650 & 3,012 \\
14 & 1,345 & 1,761 & 2,145 & 2,624 & 2,977 \\
15 & 1,341 & 1,753 & 2,131 & 2,602 & 2,947 \\
16 & 1,337 & 1,746 & 2,120 & 2,583 & 2,921 \\
17 & 1,333 & 1,740 & 2,110 & 2,567 & 2,898 \\
18 & 1,330 & 1,734 & 2,101 & 2,552 & 2,878 \\
19 & 1,328 & 1,729 & 2,093 & 2,539 & 2,861 \\
20 & 1,325 & 1,725 & 2,086 & 2,528 & 2,845 \\
21 & 1,323 & 1,721 & 2,080 & 2,518 & 2,831 \\
22 & 1,321 & 1,717 & 2,074 & 2,508 & 2,819 \\
23 & 1,319 & 1,714 & 2,069 & 2,500 & 2,807 \\
24 & 1,318 & 1,711 & 2,064 & 2,492 & 2,797 \\
25 & 1,316 & 1,708 & 2,060 & 2,485 & 2,787 \\
26 & 1,315 & 1,706 & 2,056 & 2,479 & 2,779 \\
27 & 1,314 & 1,703 & 2,052 & 2,473 & 2,771 \\
28 & 1,313 & 1,701 & 2,048 & 2,467 & 2,763 \\
29 & 1,311 & 1,699 & 2,045 & 2,462 & 2,756 \\
30 & 1,310 & 1,697 & 2,042 & 2,457 & 2,750 \\
40 & 1,303 & 1,684 & 2,021 & 2,423 & 2,704 \\
60 & 1,296 & 1,671 & 2,000 & 2,390 & 2,660 \\
120 & 1,289 & 1,658 & 1,980 & 2,358 & 2,617 \\
$\infty$ & 1,282 & 1,645 & 1,960 & 2,326 & 2,576 \\

\bottomrule
\end{tabular}
}
\end{center}

\textit{Okuma örneği:} $n=25$ için $\nu=n-1=24$ olur. Yüzde 95 çift yönlü
güven aralığında $q=0{,}975$ sütunundan $t_{0{,}975,24}=2{,}064$ okunur.
$\nu\to\infty$ satırındaki değerler standart normal kritik değerlerine yaklaşır.


\clearpage
\section{Ki-Kare Dağılımı Kritik Değerleri}
\label{sec:chi-square-table}
\textit{Kullanım:} Hücreler $\P(\chi^2_\nu\leq\chi^2_{q,\nu})=q$ olacak
birikimli olasılık nicelikleridir. Sağ kuyrukta anlamlılık düzeyi $\alpha$
için $q=1-\alpha$ alınır. Hesaplanan değer ilgili serbestlik derecesi
satırındaki kritik değerle karşılaştırılır.

\begin{center}
\scriptsize
\renewcommand{\arraystretch}{0.84}
\setlength{\tabcolsep}{3.5pt}
\resizebox{0.88\textwidth}{!}{%
\begin{tabular}{c*{5}{r}}
\toprule
$\nu$ & $q=0{,}90$ & $q=0{,}95$ & $q=0{,}975$ & $q=0{,}99$ & $q=0{,}995$ \\
\midrule
1 & 2,706 & 3,841 & 5,024 & 6,635 & 7,879 \\
2 & 4,605 & 5,991 & 7,378 & 9,210 & 10,597 \\
3 & 6,251 & 7,815 & 9,348 & 11,345 & 12,838 \\
4 & 7,779 & 9,488 & 11,143 & 13,277 & 14,860 \\
5 & 9,236 & 11,070 & 12,833 & 15,086 & 16,750 \\
6 & 10,645 & 12,592 & 14,449 & 16,812 & 18,548 \\
7 & 12,017 & 14,067 & 16,013 & 18,475 & 20,278 \\
8 & 13,362 & 15,507 & 17,535 & 20,090 & 21,955 \\
9 & 14,684 & 16,919 & 19,023 & 21,666 & 23,589 \\
10 & 15,987 & 18,307 & 20,483 & 23,209 & 25,188 \\
11 & 17,275 & 19,675 & 21,920 & 24,725 & 26,757 \\
12 & 18,549 & 21,026 & 23,337 & 26,217 & 28,300 \\
13 & 19,812 & 22,362 & 24,736 & 27,688 & 29,819 \\
14 & 21,064 & 23,685 & 26,119 & 29,141 & 31,319 \\
15 & 22,307 & 24,996 & 27,488 & 30,578 & 32,801 \\
16 & 23,542 & 26,296 & 28,845 & 32,000 & 34,267 \\
17 & 24,769 & 27,587 & 30,191 & 33,409 & 35,718 \\
18 & 25,989 & 28,869 & 31,526 & 34,805 & 37,156 \\
19 & 27,204 & 30,144 & 32,852 & 36,191 & 38,582 \\
20 & 28,412 & 31,410 & 34,170 & 37,566 & 39,997 \\
21 & 29,615 & 32,671 & 35,479 & 38,932 & 41,401 \\
22 & 30,813 & 33,924 & 36,781 & 40,289 & 42,796 \\
23 & 32,007 & 35,172 & 38,076 & 41,638 & 44,181 \\
24 & 33,196 & 36,415 & 39,364 & 42,980 & 45,559 \\
25 & 34,382 & 37,652 & 40,646 & 44,314 & 46,928 \\
26 & 35,563 & 38,885 & 41,923 & 45,642 & 48,290 \\
27 & 36,741 & 40,113 & 43,195 & 46,963 & 49,645 \\
28 & 37,916 & 41,337 & 44,461 & 48,278 & 50,993 \\
29 & 39,087 & 42,557 & 45,722 & 49,588 & 52,336 \\
30 & 40,256 & 43,773 & 46,979 & 50,892 & 53,672 \\
40 & 51,805 & 55,758 & 59,342 & 63,691 & 66,766 \\
50 & 63,167 & 67,505 & 71,420 & 76,154 & 79,490 \\
60 & 74,397 & 79,082 & 83,298 & 88,379 & 91,952 \\
80 & 96,578 & 101,879 & 106,629 & 112,329 & 116,321 \\
100 & 118,498 & 124,342 & 129,561 & 135,807 & 140,169 \\
\bottomrule
\end{tabular}
}
\end{center}

\textit{Okuma örneği:} $3\times4$ çapraz tabloda
$\nu=(3-1)(4-1)=6$'dır. Sağ kuyrukta $\alpha=0{,}05$ için $q=0{,}95$
sütunundan kritik değer $\chi^2_{0{,}95,6}=12{,}592$ bulunur.

\ifimoders\else
\clearpage
\section[F Dağılımı Kritik Değerleri]{$F$ Dağılımı Kritik Değerleri}
\label{sec:f-table}

$F_{q;\nu_1,\nu_2}$ niceliğinin solunda $q$ olasılık bulunur. Sütunlar pay
serbestlik derecesini $\nu_1$, satırlar payda serbestlik derecesini $\nu_2$
gösterir. Tek yönlü ANOVA'da $\nu_1=k-1$ ve $\nu_2=N-k$'dır; test yalnız sağ
kuyruğu kullandığından $\alpha=0{,}05$ için $q=0{,}95$ tablosu seçilir.

\begin{table}[H]
\centering
\caption{$F_{0{,}95;\nu_1,\nu_2}$ kritik değerleri: sağ kuyrukta $\alpha=0{,}05$.}
\label{tab:f-critical-05}
\small
\begin{tabular}{crrrrrr}
\toprule
$\nu_2\backslash\nu_1$ & 1 & 2 & 3 & 4 & 5 & 6 \\
\midrule
5   & 6,608 & 5,786 & 5,409 & 5,192 & 5,050 & 4,950 \\
10  & 4,965 & 4,103 & 3,708 & 3,478 & 3,326 & 3,217 \\
15  & 4,543 & 3,682 & 3,287 & 3,056 & 2,901 & 2,790 \\
20  & 4,351 & 3,493 & 3,098 & 2,866 & 2,711 & 2,599 \\
30  & 4,171 & 3,316 & 2,922 & 2,690 & 2,534 & 2,421 \\
60  & 4,001 & 3,150 & 2,758 & 2,525 & 2,368 & 2,254 \\
120 & 3,920 & 3,072 & 2,680 & 2,447 & 2,290 & 2,175 \\
$\infty$ & 3,841 & 2,996 & 2,605 & 2,372 & 2,214 & 2,099 \\
\bottomrule
\end{tabular}
\end{table}

\begin{table}[H]
\centering
\caption{$F_{0{,}99;\nu_1,\nu_2}$ kritik değerleri: sağ kuyrukta $\alpha=0{,}01$.}
\label{tab:f-critical-01}
\small
\begin{tabular}{crrrrrr}
\toprule
$\nu_2\backslash\nu_1$ & 1 & 2 & 3 & 4 & 5 & 6 \\
\midrule
5   & 16,258 & 13,274 & 12,060 & 11,392 & 10,967 & 10,672 \\
10  & 10,044 & 7,559 & 6,552 & 5,994 & 5,636 & 5,386 \\
15  & 8,683 & 6,359 & 5,417 & 4,893 & 4,556 & 4,318 \\
20  & 8,096 & 5,849 & 4,938 & 4,431 & 4,103 & 3,871 \\
30  & 7,562 & 5,390 & 4,510 & 4,018 & 3,699 & 3,473 \\
60  & 7,077 & 4,977 & 4,126 & 3,649 & 3,339 & 3,119 \\
120 & 6,851 & 4,787 & 3,949 & 3,480 & 3,174 & 2,956 \\
$\infty$ & 6,635 & 4,605 & 3,782 & 3,319 & 3,017 & 2,802 \\
\bottomrule
\end{tabular}
\end{table}

\textit{Okuma örneği:} Dört grup ve toplam $N=24$ gözlem için
$\nu_1=4-1=3$, $\nu_2=24-4=20$ olur. Yüzde 5 düzeyinde
Tablo~\ref{tab:f-critical-05}'ten $F_{0{,}95;3,20}=3{,}098$ okunur.
\fi

\section{Yazılımla tam nicelik ve kuyruk olasılığı}

Tablolar el hesabını ve kavramsal okumayı destekler; tam raporlama için
yazılımın verdiği nicelik ve \pval-değeri kullanılmalıdır.

\begin{softwarebox}
\ifimoders
\begin{lstlisting}
from scipy import stats

print(stats.norm.ppf(.975))
print(stats.t.ppf(.975, df=24))
print(stats.chi2.ppf(.95, df=6))

# Gozlenen t=2.5 icin cift yonlu p-degeri
print(2 * stats.t.sf(abs(2.5), df=24))
\end{lstlisting}
\else
\begin{lstlisting}
from scipy import stats

print(stats.norm.ppf(.975))
print(stats.t.ppf(.975, df=24))
print(stats.chi2.ppf(.95, df=6))
print(stats.f.ppf(.95, dfn=3, dfd=20))

# Gozlenen F=4.2 icin sag kuyruk p-degeri
print(stats.f.sf(4.2, dfn=3, dfd=20))
\end{lstlisting}
\fi
\end{softwarebox}